% mage-auto-heading: alignment-principle
% Authoritative LaTeX manuscript; Markdown is a generated compatibility projection.

% <!-- part-title: Alignment -->

% <!-- chapter-title: Summary -->

% <!-- coda: true -->

\chapter{Summary}

\begin{enumerate}
\def\labelenumi{\arabic{enumi}.}
\tightlist
\item
  \textbf{Enforcement is what makes an obligation more than a request.} Guidance can influence a reasoner; enforcement determines what the engineered environment will permit or accept.
\item
  \textbf{Place the mechanism where the obligation becomes decidable.} Constraints restrict actions; sensors produce evidence; gates evaluate whether work may proceed and enforce that decision. Prefer the earliest boundary where the environment has enough meaning and evidence to decide the obligation adequately.
\item
  \textbf{Correspondence is not correctness.} Models, implementations, and evidence can agree while being wrong together. Independent obligations and evidence are needed to distinguish faithful realization from mutually consistent error.
\item
  \textbf{Convert recurring judgment into durable structure.} When failures or repeated engineering decisions reveal knowledge future work should inherit, encode that knowledge into models, mechanisms, architecture, procedures, or evidence. When later work benefits from that structure, it becomes engineering capital.
\item
  \textbf{Govern the governance machinery.} Constraints, sensors, and gates form an engineering system of their own. Their coverage, correspondence, cost, effects, and failure modes must themselves be understood and controlled.
\end{enumerate}

% <!-- principlebox -->

% <!-- box-family: canonical -->

\begin{quote}
\mbox{}%
\subsection{Alignment principle}\label{heading:alignment-principle}

Make engineering obligations enforceable through mechanisms that constrain actions or states, produce evidence, or evaluate evidence to control whether work may proceed.
\end{quote}
